\section{Results}
\label{sec:results}
We distinguish published WORKS endpoints, reanalysis of retained component
measurements, and the new live executions described in
Section~\ref{sec:live-method}.

\begin{figure}[!t]
\centering
\includegraphics[width=\columnwidth]{\campReadyPath result_assets/live_margins.pdf}
\caption{CAMP request margins in the new live executions: Warm-Transfer cold, warm, and new input after warm; Direct-Transfer new input directly after the shared cold checkpoint. Margins exclude the 128~MiB execution allowance. Negative margins denote nominal requests below task peak, not observed OOMs.}
\label{fig:bowtie}
\end{figure}

\subsection{Memory-mapped access and feature contributions}
Table~\ref{tab:mmap} details the read-only versus Read + mmap comparison
briefly reported in WORKS~\cite{camp2026works}. Adding the mmap signal
reduces underallocations from 32 to 30 and total requests by 18.11\%.
Median prediction error nevertheless increases from 22.04\% to 26.51\%.
Thus, a better reservation tradeoff need not coincide with lower median
prediction error; the two-task difference alone is not evidence of a
robust advantage across datasets or random splits.

\begin{table}[!t]
\caption{Mmap ablation on the same 5,026 tasks. Native development-selected requests; no request scaling.}
\label{tab:mmap}
\centering\small
\input{\campReadyPath result_assets/mmap_table}
\end{table}

WORKS reports 87 underallocations for A+P and 31 for CAMP, with total
requests of 26,067.45 and 34,277.55~GiB~\cite{camp2026works}. The retained
reproduction gives 85 and 30; these are distinct request realizations,
not replacements for the published numbers. Table~\ref{tab:features}
expands the four feature combinations using configuration controls.
Under global prediction, full CAMP gives 74 underallocations versus 85
for A+P, with 2.06\% more requested memory. Under process-tail prediction,
A+P attains fewer underallocations (24 versus 30) but requests
603,467 versus 34,271~GiB. The full feature set therefore does not dominate
every safety metric. Relative to consumption history alone under
process-tail prediction, adding $\hat c$ reduces underallocations from 46
to 30 and requests by 57.61\%, while MdAPE increases from 21.32\% to
26.51\%. These comparisons support evaluating feature value jointly with
reservation cost.

\begin{table}[!t]
\caption{Four feature variants with configuration held fixed (5,026 tasks). U: underallocation count; R: aggregate requests in thousands of GiB.}
\label{tab:features}
\centering\small
\input{\campReadyPath result_assets/feature_table}
\end{table}

\subsection{Selective auditing and estimated net saving}
At the primary 20\% budget, the retained downstream replay gives 265.8
versus 305.8 mean underallocations for selective and random auditing:
13.1\% fewer, with 2.4\% more requested memory. This preliminary result
was also summarized in WORKS~\cite{camp2026works}. Figure~\ref{fig:audit}(a)
shows that the direction varies with budget. At 20\%, the paired 95\%
interval for the rate difference is $[-0.416,0.007]$ percentage points,
which includes zero. The five sampling repetitions do not establish
uniform improvement or robustness across model seeds.



The fixed-allocation saving is 23.596~GiB-hours before audit cost.
Figure~\ref{fig:audit}(b) gives net estimates of 17.597, 11.161, and
5.913~GiB-hours at 5, 10, and 15\% budgets. At 20\%, the estimate is
$-0.429$ with interval $[-1.077,0.346]$~GiB-hours. At 25, 30, 50, and
100\%, net estimates are $-6.087$, $-12.259$, $-42.938$, and
$-116.084$~GiB-hours. These values identify when the fixed saving covers
measured audit cost; they are not end-to-end deployment savings.


\subsection{Repeated-input learning and new-input adaptation}
Table~\ref{tab:bowtie} reports the main live experiment. Cold-to-warm
execution reduces aggregate CAMP requests by 75.34\%, from 977.79 to
241.12~GiB, while peak totals remain approximately 229.3~GiB.
Unused memory-time falls from 32.86089 to 0.01847~GiB-hours.
Warm requests cover measured peaks for 99.58\% of tasks, and median
prediction error is 1.13\%. These observations indicate more economical
requests after repeated-input feedback, conditional on the deployment
and accounting boundaries in Section~\ref{sec:live-method}.

\begin{table}[!t]
\caption{Main live experiment (Warm-Transfer): 6,000 successful tasks per stage. Requests, coverage, and unused memory exclude the execution allowance.}
\label{tab:bowtie}
\centering\footnotesize
\setlength{\tabcolsep}{3pt}
\input{\campReadyPath result_assets/warm_transfer_table}
\end{table}

On new input, peak totals rise to 273.44~GiB and requests to 319.22~GiB,
increases of 19.21\% and 32.39\% over warm execution, respectively.
Coverage falls to 94.15\%, while unused memory-time rises to
0.12805~GiB-hours. The new study therefore exposes undercoverage that is
not apparent from the repeated-input result, even though learning continues.
Alignment is particularly affected: its median prediction error increases
from 0.91\% to 15.60\%. This process-level result is obscured by the overall
new-input median of 4.73\%, since indexing and flagstat constitute most
tasks and retain different memory behavior. Figure~\ref{fig:bowtie} shows
the corresponding changes in request margins.

\subsection{Effect of the additional warm pass}
Table~\ref{tab:warm-ablation} compares the new-input branches from the
shared cold checkpoint. Without the extra warm pass, CAMP requests
343.56~GiB rather than 319.22~GiB, an increase of 7.62\%.
Peak totals differ by 1.65\%, whereas unused memory-time increases by
72.65\%, from 0.12805 to 0.22109~GiB-hours. Coverage nevertheless improves
from 94.15\% to 95.78\%, and overall median error falls from 4.73\% to
2.69\%. Alignment median error moves in the other direction, from
15.60\% to 18.96\%; an overall median alone does not establish improvement
for every process.

\begin{table}[!t]
\caption{Shared-cold-checkpoint ablation on the same 6,000 new-input tasks: Warm-Transfer includes an extra original-data warm pass; Direct-Transfer proceeds directly from cold.}
\label{tab:warm-ablation}
\centering\footnotesize
\input{\campReadyPath result_assets/direct_transfer_table}
\end{table}

The extra warm pass thus yields lower requested capacity and unused
memory-time in this execution, but lower new-input coverage. It does not
dominate the direct branch on every metric. Both new-input stages complete
all tasks with no recorded OOMs and no duplicate successful observations.
All 1,974 flagstat summaries match exactly between the new-input branches;
this checks summary agreement, not byte identity of the BAM files.
The separate execution allowance prevents interpreting this as evidence
of OOM-free operation using nominal CAMP requests alone. Because the
branches continue learning and the control node differs, this is neither
a frozen-model generalization test nor an estimate of a universal causal
benefit from additional warm training. Prediction latency and full
retraining resource costs are not isolated; no end-to-end net-saving claim
follows from these logical-task memory-time measures.
